Los resultados experimentales validan las cotas asintóticas teóricas y demuestran el impacto crítico de los factores constantes y la jerarquía de memoria en el desempeño real. En ordenamiento, se confirmó que los algoritmos evaluados pertenecen a la clase $\Theta(n \log n)$, donde \textit{std::sort} y \textit{Quick Sort} superan sistemáticamente a \textit{Merge Sort} gracias a su operación \textit{in-place} y óptima localidad espacial. Asimismo, se evidenció la extrema sensibilidad de \textit{Patience Sort}, alcanzando su mejor caso $\Omega(n)$ ante secuencias decrecientes al operar con una sola pila ($k=1$), pero degradándose drásticamente en arreglos ascendentes y dominios amplios ($D7$) debido a la sobrecarga del min-heap ($k=n$). En multiplicación de matrices, los datos reflejan que, si bien \textit{Strassen} es asintóticamente superior ($o(n^3)$), la sobrecarga de 18 operaciones intermedias y la asignación dinámica de memoria favorecen al algoritmo iterativo tradicional en dimensiones moderadas, alcanzando el umbral de cruce práctico recién en $n = 1024$, donde la reducción en la complejidad temporal comienza a compensar su mayor consumo espacial.